\paragraph{The remaining endpoint recipes.}
For the conditional coefficient, use the mean and binary-cubic bounds
from $V=-m/q,m/p$ and $m^2=(1-\theta)pq$. Explicitly, if
$p\in[l,u]\subset(0,1/2)$ and $\theta\le\Theta$, the positive mean
lies between $\sqrt{(1-\Theta)(1-u)/u}$ and $\sqrt{(1-l)/l}$,
the negative magnitude mean between $\sqrt{(1-\Theta)l/(1-l)}$ and
$\sqrt{u/(1-u)}$, and
\[
 \sqrt{(1-\Theta)^3(1-2u)^2/[u(1-u)]}
 \le\gamma_V\le(1-2l)/\sqrt{l(1-l)}.
\]
In $F$, the edge stationary points are
$a=t/2-x/3$ and $x=3t/4-a$, clipped to their edges; the interior
point is $x=a=3t/8$. Retain the actual common $\rho$ interval.
Its lower bound may be raised to $\rho_B(u)$ only if
$3\min(A_l,B_l)\ge(3/2)\sup_I\rho_B$; the conditional cubic identity
and $(1-\theta)^{3/2}\ge1-3\theta/2$ then prove that license.
Intersect the signed $K_1$ interval with its previously proved error interval.

For the central constant, the cheaper alternative to the complete Bernstein
rectangle follows from
\[
 G_B=B(A+B)(\rho-B),\quad
 G(A,\rho;\rho)=\rho^4/12+\rho^3A/6-\rho A^3/6+A^4/12.
\]
The latter has derivatives $(\rho-A)^2(\rho+2A)/6$ in $A$ and
$\rho^2+\rho A$ twice in $\rho$, so use the upper $A$ endpoint and
both $\rho$ endpoints. Intersect the weighted outer budget with
$T_+\le b_{+,u}-m_lB_l^2$ and $T_-\le b_{-,u}-m_lA_l^2$.
For $U=J_1+(\mu_4-3)/6$, add the signed central and outer intervals
and the fourth-moment interval. Use the complete $K_1,U,\phi$ rectangle
in $|CK_1+\phi U|$; do not separate the two absolute values.

\begin{lemma}[Complete central rectangle]\label{j:local:entry201:central}
For $G=\rho P_1-P_2$, let $b_{ij}(\rho)$ be its degree-four tensor
Bernstein coefficients after mapping the proved $(A,B)$ rectangle to
$[0,1]^2$. Here one may take
\[
 A_u=\sqrt{u/(1-u)},\quad B_u=\sqrt{(1-l)/l},\quad
 m_l=\sqrt{(1-\Theta)\min_{p=l,u}pq},\quad
 m_u=\sqrt{\max_{p\in[l,u]}pq}.
\]
Then
\[
 \rho J_{1,c}-J_{2,c}\le
 \max_{0\le i,j\le4}\{m_l b_{ij}(\rho),m_u b_{ij}(\rho)\}.
\]
This upper bound is convex in $\rho$.
\end{lemma}
The preceding positive-basis argument proves the lemma, including closed
faces and degenerate rectangles. The fixed algebraic objects record the
seven moment matrices, the displayed derivative and square identities,
and all25 inverse basis identities. Exact coefficient comparison checks
these finite identities; the integral and common-budget arguments above
supply their mathematical meaning. No acquisition or Fourier target
follows from coefficient acceptance alone.
